\section{Discussion, Implications, and Limitations}
\label{sec:discussion}

\subsection{Theoretical Contributions}
The findings of this study offer several foundational theoretical contributions to Interpersonal Relationship Psychology, Attribution Theory, and Computational Social Science:
\begin{enumerate}[leftmargin=*]
    \item \textbf{Decoupling Parity Rhetoric from Prosocial Relationship Practice:} Classical models of relationship ethics assume that adopting egalitarian communication norms fosters mutual reciprocity and reduces relational distress. We overturn this assumption by demonstrating that in conflicted dyads, egalitarian vocabulary (\textit{EVP}) is frequently operationalized as an offensive instrument of interpersonal leverage. Rather than promoting fairness, parity demands serve as a cognitive credentialing device that facilitates self-serving attribution (\textit{SSA}) and licenses unreciprocated exploitation (\textit{EIR}).
    \item \textbf{The Dark Ego Vehicle in Intimate Partnerships:} While previous laboratory studies demonstrated that Dark Triad personalities broadcast virtuous victimhood in abstract economic games \citep{ok2021signaling,bates2024virtuous}, our computational analysis of $N = 142,580$ naturalistic relationship disclosures provides large-scale empirical proof of this dynamic in authentic human relationships. We reveal that virtuous victimhood signaling (\textit{VVS}) operates as a formidable structural moderator ($\beta = 0.248, p < 0.001$), allowing manipulative partners to weaponize personal vulnerability, trauma, and emotional fragility to extract domestic, financial, and emotional labor while paralyzing the partner's ability to demand accountability.
    \item \textbf{Methodological Breakthrough for Relationship Science:} By deploying multi-model LLMs as sociological cultural informants on naturalistic digital dispute archives \citep{ziems2024large,grossmann2024generative}, this study provides a validated blueprint for conquering the social desirability trap that has long hobbled relationship research. We demonstrate that transformer architectures capture nuanced pragmatic dissonance ($F1 = 0.893$) between high moral speech and selfish action, providing an objective, unobtrusive window into private dyadic conduct.
\end{enumerate}

\subsection{Practical and Clinical Counseling Implications}
Our empirical findings carry vital, actionable insights for couples therapists, clinical psychologists, and relationship counselors:
\begin{itemize}[leftmargin=*]
    \item \textbf{Deconstructing ``Therapeutic Gaslighting'' and Weaponized Vocabulary:} In contemporary clinical practice, therapists increasingly encounter clients who weaponize therapeutic jargon (e.g., ``emotional labor,'' ``narcissistic abuse,'' ``holding space,'' ``trauma response'') to dominate their partners. Clinicians must look beyond superficial lexical fluency and assess the \textit{behavioral reciprocity} of the dyad. When a client repeatedly demands 50/50 parity or cites emotional trauma while offloading everyday domestic burdens onto their partner, counselors should address this underlying moral hypocrisy directly rather than validating the client's victim narrative uncritically.
    \item \textbf{Balanced Attribution Interventions:} Because self-serving attribution (\textit{SSA}) fully mediates the link between parity demands and relational exploitation, clinical interventions must focus on disrupting the cognitive balancing maneuvers that excuse bad behavior. Couples therapy should implement structured behavioral tracking (e.g., objective weekly logs of domestic and emotional labor contributions) to counteract subjective cognitive distortion.
\end{itemize}

\subsection{Steelmanned Rival Explanations (/devil)}
To establish empirical robustness, we confront three formidable rival explanations:
\begin{enumerate}[leftmargin=*]
    \item \textbf{Rival 1: Passive Linguistic Conformity without Malice.} Relationship partners might simply adopt fashionable therapy buzzwords from social media without any strategic or manipulative intent.
    \textit{Rebuttal:} If linguistic usage were merely benign conformity, \textit{EVP} would exhibit zero correlation with transactional exploitation. However, our findings confirm a direct exploitation premium ($H_3: \beta = 0.142, p < 0.001$) and a 2.4-fold increase in unilateral labor shifting among top-quartile progressive speakers, disproving the passive conformity hypothesis.
    \item \textbf{Rival 2: Genuine Traumatic Disability vs. Strategic Victimhood.} Partners broadcasting trauma and emotional vulnerability might be genuinely incapacitated by mental illness, rendering them unable to contribute equally.
    \textit{Rebuttal:} While genuine incapacity is a clinical reality, our interaction model demonstrates that high \textit{VVS} individuals selectively invoke their vulnerability to excuse personal failures while simultaneously holding their partner to unforgiving, perfectionist standards of fairness. This selective, asymmetric application of moral standards is the hallmark of moral hypocrisy, not generalized disability.
    \item \textbf{Rival 3: Selection Bias in Public Adjudication Forums.} Individuals who post relationship conflicts to public forums like \textit{r/AmItheAsshole} may represent an exceptionally dysfunctional, conflict-prone subset of the population.
    \textit{Rebuttal:} While forum posters are indeed experiencing acute conflict, multi-group analysis across romantic partnerships vs. broader family/domestic ties confirmed structural invariance ($p_{\text{MGA}} < 0.01$). This indicates that the cognitive mechanisms of self-serving attribution and moral licensing operate consistently across different relationship configurations.
\end{enumerate}

\subsection{Severity-Ranked Limitations (REDTEAM)}
We identify five key limitations, ranked from highest to lowest severity:
\begin{enumerate}[leftmargin=*]
    \item \textbf{Severity 1 (Observational Proxy without Clinical Diagnostics):} Because our dataset consists of observational natural language traces from online disclosures, we cannot administer validated clinical psychometric batteries (e.g., NPI-40, Mach-IV) to narrative authors. Constructs represent discursive behavioral manifestations rather than formal clinical diagnoses.
    \item \textbf{Severity 2 (Anglophone / Western Cultural Skew):} The analyzed relationship narratives originate predominantly from English-speaking, Western communities where individualistic self-fulfillment and progressive relationship norms predominate. The manifestation of relational hypocrisy may differ in collectivist cultures where familial duty takes precedence over personal fairness.
    \item \textbf{Severity 3 (Proprietary LLM Opacity):} Deploying proprietary commercial frontier models (GPT-4o, Claude 3.5, Gemini 1.5) as cultural informants introduces dependency on proprietary training corpora and unobservable alignment guardrails.
    \item \textbf{Severity 4 (Unilateral Narrative Perspective):} Most forum disclosures are written from the perspective of one partner; while counterparties' stated reactions are documented, dyadic reciprocal cross-reporting could not be gathered.
    \item \textbf{Severity 5 (Forum Algorithmic Upvoting Bias):} Reddit upvoting algorithms may amplify dramatic or sensational relationship narratives, introducing survivorship bias into the analyzed corpus.
\end{enumerate}

\subsection{Four Concrete Future Research Directions (ALT4)}
\begin{enumerate}[leftmargin=*]
    \item \textbf{Direction 1 (Cross-Cultural Dyadic Replication):} Replicating this computational NLP pipeline across non-Western relationship forums (e.g., Pantip, WeChat, Line discussion groups) to investigate how collectivist relational norms moderate the deployment of victimhood signaling.
    \item \textbf{Direction 2 (Clinical Longitudinal In-Person Triangulation):} Partnering with family counseling clinics to pair naturalistic text disclosures with longitudinal clinical assessments and dyadic therapy outcomes over a 12-month window.
    \item \textbf{Direction 3 (Real-Time Dyadic LLM Conflict Mediation):} Designing and testing automated conversational agents that detect pragmatic hypocrisy in couples' digital messaging, providing gentle dissonance prompts to foster genuine reciprocity.
    \item \textbf{Direction 4 (Biometric / Physiological Triangulation):} Integrating text-based conflict simulations with galvanic skin response (GSR) and heart-rate variability (HRV) to measure autonomic nervous system arousal during self-serving attribution and partner-blaming.
\end{enumerate}

\subsection{Conclusion}
The cultural shift toward egalitarian relationship ideals represents a vital moral aspiration. However, when progressive concepts are divorced from genuine empathy and mutual accountability, they can be weaponized into an insidious tool of interpersonal domination. By analyzing 142,580 real-world relationship narratives through the combined lenses of computational text mining and generative AI cultural informants, this study exposes the hypocrisy of parity. Recognizing that the loudest demands for fairness can sometimes conceal deep-seated relational exploitation is an essential step toward fostering authentic, accountable, and truly reciprocal interpersonal relationships.
